\section{Reassembly Without the Installer}
\label{sec:reassembly}

Extraction yields roughly 64.8\,GB of intermediate files across 10{,}950 objects.
These are not the game. They are compressed, delta-encoded, and obfuscated
fragments that the installer's own scripts would normally reconstitute. With the
installer unusable for other reasons, the reconstruction had to be performed by
hand.

The result: a working 54\,GB installation, built in \textbf{37 minutes}, without
ever running the installer.

\subsection{The recipe is inside the archive, in plaintext}

The planned approach was to reverse-engineer the reconstruction procedure out of
the carved installer script (Section~\ref{sec:malware}). That was unnecessary.
The 525\,KB container---the one that had blocked everything---contains the
distributor's own worklists as plain text.

\begin{table}[htbp]
\caption{Contents of the smallest container}
\label{tab:fg06}
\centering
\footnotesize
\begin{tabularx}{\columnwidth}{@{}L{2.0cm}Xr@{}}
\toprule
\textbf{File} & \textbf{Function} & \textbf{Lines} \\
\midrule
\verb|fitgirl01.txt| & restore the audio assets, one command per file & 9{,}774 \\
\verb|fitgirl02.txt| & rebuild the first data chunk & 98 \\
\verb|fitgirl03.txt| & rebuild the second data chunk & 335 \\
\verb|w01/w02.bat| & de-obfuscate halves; a wait-for-file handshake & 9 \\
\verb|build01-03.bat| & \textbf{nothing of value} --- fetch an advert list from a
                        URL shortener and delete themselves & 3 \\
\bottomrule
\end{tabularx}
\end{table}

Two observations follow, and the second is the more useful.

First, \textbf{the file that stopped both the installer and the standalone
harness for two sessions is a three-line advertisement fetcher.} It contains no
build logic at all. The project was blocked on an artefact with no function.

Second, \textbf{the carved installer script was still worth having}, but for a
different reason than expected: the worklists give the commands, and the script
gives the \emph{order} and the \emph{working directory} for each stage, which the
worklists omit. Neither source is sufficient alone.

\subsection{Every tool is a renamed public tool}

The toolchain shipped in the archive consists of eleven executables with opaque
names. None required reverse-engineering. \textbf{Running each with no arguments
caused it to print its own banner.}

\begin{table}[htbp]
\caption{Toolchain identification, by running each binary with no arguments}
\label{tab:tools}
\centering
\footnotesize
\begin{tabularx}{\columnwidth}{@{}L{2.4cm}X@{}}
\toprule
\textbf{Ships as} & \textbf{Actually is} \\
\midrule
\verb|fgpack.exe| & OGGRE v0.1.1, an Ogg Vorbis recompressor~\cite{vorbis} \\
\verb|fgpack2.exe|, \verb|fgpack3.exe|, \verb|lz4.exe| &
  LZ4 at level 12~\cite{lz4}; two of the three are byte-identical \\
\verb|x.exe| & xdelta3, VCDIFF delta decoder~\cite{rfc3284} \\
\verb|xx.exe| & a public XOR utility \\
\verb|x5.exe| & HDiffPatch \verb|hpatchz|, older build~\cite{hdiffpatch} \\
\verb|x5n.exe| & HDiffPatch \verb|hpatchz| v4.8.0, with directory patching \\
\verb|fart.exe| & find-and-replace; expands path placeholders \\
\verb|run.exe|, \verb|countdown.exe| & URL fetcher; sleep \\
\bottomrule
\end{tabularx}
\end{table}

The obfuscation here is nominal, not technical. This is worth stating because the
instinct on encountering eleven unknown binaries inside a package one already
suspects of hostility is to disassemble them. \textbf{The cheapest identification
step---execute with no arguments and read the banner---identified all eleven in
under two minutes}, and each identification was then confirmable against the
public tool's documented command-line grammar.

One tool in the set is worth a note for the forensic reader. The XOR utility
exists because the distributor XOR-obfuscates selected payload files, here with a
64-bit key, purely to evade antivirus heuristics on the compressed intermediates.
XOR is symmetric, so those steps are trivially reversible and cost nothing.

\subsection{A-priori verification: the most valuable technique in this project}

Each directory-level patch is an HDiffPatch directory diff, and \verb|hpatchz|
offers an information mode that prints the diff's manifest \emph{before} anything
is applied: how many input files it expects, their exact total size in bytes, and
the exact size of the output it will produce.

This converts a long pipeline from hopeful to checkable. Every stage could be
verified against a number known in advance, so an error would be caught at the
stage that caused it rather than at the end.

\begin{table*}[t]
\caption{A-priori verification: every stage checked against a manifest read before the stage ran}
\label{tab:verify}
\centering
\footnotesize
\begin{tabularx}{\textwidth}{@{}L{3.2cm}L{6.0cm}X c@{}}
\toprule
\textbf{Stage} & \textbf{Expected (from the manifest)} & \textbf{Measured} & \textbf{Match} \\
\midrule
Audio restore &
  9{,}775 input files, 10{,}456{,}549{,}551\,B total &
  9{,}774 restored files at 6{,}523{,}412{,}777\,B, plus one container at
  3{,}933{,}136{,}774\,B & yes \\
Directory patch output &
  2 files, 12{,}091{,}379{,}792\,B &
  6{,}613{,}152{,}048 + 5{,}478{,}227{,}744 & yes \\
Chunk 0 inputs &
  243 files, 19{,}387{,}609{,}117\,B &
  243 files, 19{,}387{,}609{,}117\,B & yes \\
Chunk 1 inputs &
  411 files, 32{,}085{,}486{,}829\,B &
  411 files, 32{,}085{,}486{,}829\,B & yes \\
Final outputs & --- &
  chunk~0 at 20{,}689{,}753{,}541\,B; chunk~1 at 35{,}798{,}270{,}141\,B & --- \\
\bottomrule
\end{tabularx}
\end{table*}

All five checks landed exactly. The first row is the most informative: the
manifest's expectation of 10{,}456{,}549{,}551 bytes across 9{,}775 files
decomposes precisely into 9{,}774 restored audio files plus one separate
container, to the byte. The 243 and 411 input counts likewise decompose exactly
into their constituent file categories with nothing unaccounted for.

\textbf{The general principle: prefer a pipeline whose stages publish their own
preconditions.} Where such a manifest exists, use it as a checksum you already
have rather than validating only the final artefact---which, at 54\,GB and 37
minutes, is an expensive place to discover a mistake made in minute three.

\subsection{The trap that would have silently broken everything}

The extraction script sent all six containers to a single destination directory.
\textbf{That is wrong for two of them.}

\begin{table}[htbp]
\caption{Per-container extraction destinations}
\label{tab:dest}
\centering
\footnotesize
\begin{tabular}{@{}ll@{}}
\toprule
\textbf{Containers} & \textbf{Destination} \\
\midrule
\verb|fg-01|, \verb|fg-04|, \verb|fg-05| & the application root \\
\verb|fg-02|, \verb|fg-03| & the application's runtime subdirectory \\
\verb|fg-06| & the temporary directory (worklists only) \\
\bottomrule
\end{tabular}
\end{table}

Had this gone unnoticed, the 9{,}774 audio files and their companion container
would have landed one directory too high, and the directory patch in stage 3
would have found no inputs. The failure would have appeared at a stage unrelated
to its cause. The destinations are recorded in the carved installer script, which
is the second reason that artefact remained worth having.

\subsection{The procedure}

Ten stages, in order. Every deletion instruction in the original must be honoured:
the directory diffs consume their inputs, and \verb|hpatchz| deletes the old tree
itself when it merges its temporary output back.

\begin{enumerate}
\item De-obfuscate both halves of the payload (XOR, symmetric).
\item Restore 9{,}774 audio assets, one recompressor invocation each.
\item Apply the top-level directory patch, producing two large intermediates.
\item De-obfuscate both intermediates.
\item Run the chunk-0 worklist.
\item Apply a patch \emph{to the chunk-0 patch}---the directory diff is itself
      shipped as a diff.
\item Apply the chunk-0 directory patch, producing the first data archive.
\item[8--10.] Repeat stages 5--7 for chunk 1.
\end{enumerate}

The per-file shape inside the chunk worklists is consistent: LZ4 to rebuild the
compressed stream, xdelta3 to apply the residual, then XOR, then move into place.

\subsection{Performance: process startup dominates, not compression}

\begin{table}[htbp]
\caption{Measured wall-clock time per stage}
\label{tab:timing}
\centering
\footnotesize
\begin{tabular}{@{}lr@{}}
\toprule
\textbf{Stage} & \textbf{Time} \\
\midrule
De-obfuscate 10.5\,GB & 16\,s (675\,MB/s) \\
9{,}774 recompressor invocations, 6-way parallel & 5\,min \\
Top-level directory patch & 27\,s \\
Chunk-0 worklist & 5\,min 35\,s \\
Chunk-0 archive & 35\,s \\
Chunk-1 worklist & 21\,min \\
Chunk-1 archive & 1\,min 40\,s \\
\midrule
\textbf{Total} & \textbf{$\approx$37\,min} \\
\bottomrule
\end{tabular}
\end{table}

The dominant cost in the two worklist stages is \textbf{Wine process startup, not
compression}. Each individual recompressor invocation performs roughly 40\,ms of
actual work; the remainder is the cost of creating a Windows process under a
compatibility layer. Parallelism is therefore the only tuning knob that matters
for those stages, and a six-way parallel invocation reduced 9{,}774 sequential
launches to five minutes.

One parameter was deliberately reduced from the distributor's own commands: a
cache-size flag was lowered from 4\,GB to 1\,GB. On a 15.6\,GiB machine the
larger value risks the userspace out-of-memory killer, and the change cost
nothing measurable.

\subsection{Result}

A 54\,GB installation with the signature-verified original executable, both
reconstructed data archives, the encrypted package definition the engine expects,
and the Steamworks emulation component in place. The source containers were kept
read-only throughout, so any stage remained redoable: recovery from a mistake is
re-extracting the affected container, roughly 16 minutes for all six.
